% Compile with XeLaTeX: xelatex EE180_HW1.tex (run twice).
\documentclass[11pt,a4paper]{article}

\usepackage{amsmath,amssymb}
\usepackage{geometry}
\geometry{margin=2.5cm}
\usepackage{graphicx}
\usepackage{enumitem}
\usepackage{hyperref}
\usepackage{xeCJK}

\title{{\Large\textbf{EE180 Fall Homework 1}} \\[4pt]
{\large\textbf{Due: October 8, 2026, 23:59}}}
\date{}

\begin{document}
\maketitle

\begin{center}
\textbf{Name:\quad\underline{\hspace{4cm}}\quad
ID:\quad\underline{\hspace{4cm}}}
\end{center}
\vspace{0.5em}
\noindent\rule{\textwidth}{0.4pt}

\noindent\textbf{Requirements:}
\begin{itemize}[leftmargin=1.5em]
    \item Use MATLAB or Python; Python is recommended.
    \item Submit a ZIP archive named \texttt{EE180-ID-Name-hw1.zip}, containing:
    \begin{itemize}
        \item \texttt{codes/} folder with your source code,
        \item \texttt{image/} folder with generated figures,
        \item \texttt{report-ID-Name-hw1.pdf},
        \item \texttt{predictions\_stage1.csv} and \texttt{predictions.csv} for the two stages of Question 3.
    \end{itemize}
    \item Use relative file paths and include brief instructions for running your code.
    \item Plot all time--frequency representations using magnitude in decibels:
    \[
        M_{\mathrm{dB}}(t,f)=20\log_{10}|X(t,f)|,
    \]
    where $X(t,f)$ denotes the normalized transform coefficients. Use the same
    normalization reference and color scale for plots compared within the
    same transform. For CWT, use the corresponding coefficient magnitude
    and map scales to frequency. State the normalization used, and apply a
    small positive floor before taking the logarithm to avoid $\log_{10}(0)$.
    \item Write your report in English.
\end{itemize}

\newpage
\section*{1. Fourier and Time--Frequency Analysis}

Use \textbf{Signal A} and \textbf{Signal B} in \path{data/q1_signals.mat} or
\path{data/q1_signals.npz}. Both files contain \texttt{signal\_A},
\texttt{signal\_B}, \texttt{time\_s} (seconds), and
\texttt{sampling\_rate\_hz}. The sampling frequency is $60\,\mathrm{MHz}$.
Use $\mu\mathrm{s}$ and $\mathrm{MHz}$ for your time and frequency axes.

\begin{enumerate}[label=(\alph*), leftmargin=1.5em]
    \item \textbf{Time-domain and FFT analysis.}

    Plot both signals in the time domain and their single-sided FFT magnitude
    spectra, using the same axis limits for corresponding plots and the same
    FFT normalization. Describe the observed
    differences and explain why the FFT magnitude spectra alone do not clearly
    distinguish their temporal structures.

    \item \textbf{STFT analysis.}

    Plot the STFT magnitude spectrogram of each signal using a Hann window
    of $240$ samples, a hop size of $6$ samples, and
    $N_{\mathrm{FFT}}=4096$. Plot the magnitude in dB as specified in the
    Requirements, using a common normalization reference and color scale
    for both signals. State how signal boundaries are handled.

    Describe the differences between the two signals revealed by the STFT,
    supporting your interpretation with evidence from the plots.
\end{enumerate}

\textbf{Submission:} Include your plots, explanations, and reproducible code.
\newpage
\section*{2. STFT, CWT, and Time--Frequency Resolution}

Use \textbf{Signal A} from Question 1 to investigate the effect of window
length on STFT analysis and compare the results with continuous wavelet analysis.

\begin{enumerate}[label=(\alph*), leftmargin=1.5em]
    \item \textbf{STFT analysis and signal components.}

    Choose a short window and a long window, and compute and plot the STFT
    magnitude spectrogram of Signal A for each. State and justify your
    parameter choices. Plot the magnitude in dB as specified in the Requirements.
    Use the same window type, magnitude normalization reference, and color
    scale for the short- and long-window results.

    \textit{Hint:} You may explore several window lengths, but only two representative
    results, one short and one long, are required.

    Use the STFT results together to describe the distinct signal
    components present in Signal A. Estimate their characteristic frequencies,
    occurrence times, and frequency evolution where applicable. For components
    that are close in time or frequency, estimate their separation and
    indicate which spectrogram supports each estimate.

    Support your interpretation with annotations, enlarged time--frequency
    regions, or one-dimensional slices through the STFT magnitude.

    \item \textbf{Time--frequency resolution trade-off.}

    Using examples from your results, explain the trade-off between time
    and frequency resolution as the window length changes. Discuss whether
    one window length can adequately resolve all components of Signal A.

    \item \textbf{Continuous wavelet analysis.}

    Analyze Signal A using the continuous wavelet transform with a complex
    Morlet wavelet, and plot its magnitude in dB against time and frequency,
    following the Requirements. State the wavelet parameters, scale or
    frequency range, and coefficient normalization used.

    Compare your result with the STFT results in part (a). Discuss how well
    the different signal components can be distinguished, and explain your
    observations in terms of the time and frequency localization of the
    STFT window and the wavelet. Does CWT eliminate the time--frequency resolution
    trade-off? Support your answer with evidence from your results.

    \textit{Hint:} You may explore different wavelet parameters, but presenting
    one representative CWT result is sufficient.
\end{enumerate}

\textbf{Submission:} Include your STFT and CWT plots, parameter
choices, component identification, discussion, and reproducible code.
\newpage
\section*{3. Classifying Classroom Ultrasound Recordings}

In class, you will acquire ultrasound signals from three objects or sample
configurations following the TA's demonstration. Each student will make
three independent acquisitions per class. Differences in handling may cause
recordings of the same class to vary.

Each supplied recording is a \textbf{one-dimensional signal with 896 samples}.
The \textbf{sampling frequency is $26.5216\,\mathrm{MHz}$}. The supplied signals
have already been normalized to the range $[-1,1]$; no raw-ADC normalization
is required.

You will receive:
\begin{itemize}[leftmargin=1.5em]
    \item A labeled reference set collected by the TA:
    \textbf{15 independent recordings per class} (45 in total).
    \item \textbf{Stage 1:} An additional set collected by the TA, with
    \textbf{5 recordings per class} (15 in total), separate from the reference set,
    anonymized and shuffled, with class labels withheld.
    \item \textbf{Stage 2:} The classroom recordings, anonymized and shuffled,
    with class labels withheld.
\end{itemize}

Use signal analysis methods introduced in class to extract features from the
ultrasound recordings and classify each recording in both stages as \textbf{Class 1,
2, or 3}. Submit \textbf{one label per recording}. A README in the dataset
folder will describe the file formats, identifiers, data layout, units,
and acquisition metadata.

\textbf{Signal decomposition and feature extraction.}

Apply appropriate signal decomposition or transformation methods to analyze
the reference recordings. Represent each recording by one set of features.
Compute descriptive statistics of the feature values for each class, examine
differences between classes and variation within each class, and support your
observations with plots or summary tables. Apply the same preprocessing and
feature-extraction procedure to the reference set and both task stages.
If you include time--frequency plots, follow the dB display convention in the
Requirements.

\textbf{Classification.}

Use the feature statistics and plots from the labeled reference set to
\textbf{manually choose thresholds or characteristic ranges} for distinguishing
the three classes. State your decision rules as explicit numerical conditions,
explain how you chose them, and apply the same rules in both stages.

As an \textbf{optional alternative}, you may form a feature vector for each
recording and train a machine-learning classifier using the labeled reference
set. \textbf{Machine learning is not required.}

Describe your features, decision rules or optional model, and parameter choices.
Use the labeled reference set to determine your thresholds or model parameters.
Base predictions on the recorded signals rather than acquisition order or
knowledge of which object was measured.

\textbf{Two-stage task.}

In \textbf{Stage 1}, classify the 15 additional TA recordings. In
\textbf{Stage 2}, classify each shuffled classroom recording. Labels are withheld
in both stages; submit one predicted label per recording. The TA will evaluate
both submissions against private ground-truth labels. Students are not required
to report task-stage accuracy or analyze errors against withheld labels.

\textbf{Submission and assessment:} \textbf{Both stages must be submitted and
are included in the assessment.} Submit \texttt{predictions\_stage1.csv} for
Stage 1 and \texttt{predictions.csv} for Stage 2. Each file must contain the
columns \texttt{sample\_id} (the recording identifier) and
\texttt{predicted\_class}, with one row per recording in that stage. Encode
\texttt{predicted\_class} as 1, 2, or 3. Describe your method and reference-set analysis
in the report, including supporting figures, and provide reproducible code.

\end{document}
